% --- Archon physics formalization source begin ---
% archon:physics
% archon:covers PhyXMiniProblems/problem_phyx_mini_0722.lean
% archon:source-report reports/phyx_mini/problem_phyx_mini_0722.source.json

\chapter{Physics problem phyx\_mini\_0722}
\label{ch:PhyXMiniProblems_problem_phyx_mini_0722}

\paragraph{Problem source.}
\#\# Physical scenario

A \textbackslash{}(10\textbackslash{},\textbackslash{}mathrm\{cm\} \textbackslash{}times 10\textbackslash{},\textbackslash{}mathrm\{cm\} \textbackslash{}times 10\textbackslash{},\textbackslash{}mathrm\{cm\}\textbackslash{}) block of wood with density \textbackslash{}(700\textbackslash{},\textbackslash{}mathrm\{kg/m\textasciicircum{}3\}\textbackslash{}) is held underwater by a string tied to the bottom of the container.

\#\# Question

What is the tension in the string

\#\# Answer choices

- A: A: 3.5N

- B: B: 3.2N

- C: C: 2.9N

- D: D: 3.8N

\#\# Figure caption (auxiliary; use the image as primary evidence)

The image depicts a block of wood submerged in a liquid, illustrating the concept of buoyancy. Here's a detailed description:

1. **Objects and Scene:**
   - A block of wood is shown submerged in a container of liquid.
   - The block is attached to a string at the bottom.

2. **Forces:**
   - **Buoyant Force (\textbackslash{}( \textbackslash{}vec\{F\}\_B \textbackslash{})):** An upward arrow labeled with \textbackslash{}( \textbackslash{}vec\{F\}\_B \textbackslash{}) represents the buoyant force acting on the block.
   - **Gravitational Force (\textbackslash{}( \textbackslash{}vec\{F\}\_G \textbackslash{})):** A downward arrow labeled with \textbackslash{}( \textbackslash{}vec\{F\}\_G \textbackslash{}) shows the gravitational force on the block.
   - **Tension (\textbackslash{}( \textbackslash{}vec\{T\} \textbackslash{}))**: Another downward arrow labeled \textbackslash{}( \textbackslash{}vec\{T\} \textbackslash{}) represents the tension in the string.

3. **Text:**
   - A label, "The buoyant force pushes up on the block," is positioned to the left, associated with a dotted line pointing to the upward force arrow.
   - The block is labeled "Block of wood."
   - The string is labeled "String."

The image visually represents the relationship between the forces acting on an object submerged in a fluid, showing how buoyant force counteracts gravity.

\#\# Dataset metadata

Statics; Physical Model Grounding Reasoning, Multi-Formula Reasoning

\paragraph{Recorded answer/context.}
C: C: 2.9N

\paragraph{Figure/image path.}
/root/proposal\_for\_physic/science-mango/phyx\_mini\_run/phyx\_data/test\_image/722.png

\paragraph{Formalization target.}
create a compiling Lean file with sorry bodies at `PhyXMiniProblems/problem\_phyx\_mini\_0722.lean`. The Lean declarations must preserve the physical quantities, dimensions or dimensional roles, figure labels, governing-law hypotheses, and final relation expressed by this problem.
Use Mathlib/Physlib names found through LeanExplore where available. If a physics API is missing, introduce faithful local abstractions rather than scalar placeholder aliases.

\begin{theorem}[Physics formalization target]
\label{thm:physics:phyx_mini_0722:target}
\lean{PhyXMiniProblems.ProblemPhyXMini0722.problem_phyx_mini_0722}
\uses{def:physics:phyx-mini-0722:phyxminiproblems-problemphyxmini0722-forceinnewtons, def:physics:phyx-mini-0722:phyxminiproblems-problemphyxmini0722-submergedwoodblocksetup, def:physics:phyx-mini-0722:phyxminiproblems-problemphyxmini0722-matchessubmergedwoodblockscenario, def:physics:phyx-mini-0722:phyxminiproblems-problemphyxmini0722-matchessuppliedsubmergedblockfigure, def:physics:phyx-mini-0722:phyxminiproblems-problemphyxmini0722-matchessubmergedblockproblemreadouts, def:physics:phyx-mini-0722:phyxminiproblems-problemphyxmini0722-usesstandardwaterandgravity, def:physics:phyx-mini-0722:phyxminiproblems-problemphyxmini0722-hasphysicalsubmergedblockparameters, def:physics:phyx-mini-0722:phyxminiproblems-problemphyxmini0722-satisfiescubeanddisplacementlaws, def:physics:phyx-mini-0722:phyxminiproblems-problemphyxmini0722-satisfiessubmergedblockstaticslaws, def:physics:phyx-mini-0722:phyxminiproblems-problemphyxmini0722-answerchoice, def:physics:phyx-mini-0722:phyxminiproblems-problemphyxmini0722-displayedtensioninnewtons, def:physics:phyx-mini-0722:phyxminiproblems-problemphyxmini0722-roundstonearesttenthnewton, def:physics:phyx-mini-0722:phyxminiproblems-problemphyxmini0722-isnearestdisplayedtension}
The assigned autoformalize agent should translate this physics problem into Lean declarations in the covered file, with theorem and lemma proof bodies written as `by sorry`.
\end{theorem}
\begin{proof}
This is an autoformalization task, not a proof task. Produce statements that can later be proved by the physics prover without weakening the physical model.
\end{proof}

% --- Archon declaration topology begin ---
\section{Declaration topology}
\begin{definition}[volume Dimension]
  \label{def:physics:phyx-mini-0722:phyxminiproblems-problemphyxmini0722-volumedimension}
  \lean{PhyXMiniProblems.ProblemPhyXMini0722.volumeDimension}
  The physical dimension of volume, L³
\end{definition}
\begin{proof}
  This is the declared model definition of volume Dimension.
\end{proof}
\begin{definition}[mass Density Dimension]
  \label{def:physics:phyx-mini-0722:phyxminiproblems-problemphyxmini0722-massdensitydimension}
  \lean{PhyXMiniProblems.ProblemPhyXMini0722.massDensityDimension}
  The physical dimension of mass density, M L⁻³
\end{definition}
\begin{proof}
  This is the declared model definition of mass Density Dimension.
\end{proof}
\begin{definition}[acceleration Dimension]
  \label{def:physics:phyx-mini-0722:phyxminiproblems-problemphyxmini0722-accelerationdimension}
  \lean{PhyXMiniProblems.ProblemPhyXMini0722.accelerationDimension}
  The physical dimension of acceleration, L T⁻²
\end{definition}
\begin{proof}
  This is the declared model definition of acceleration Dimension.
\end{proof}
\begin{definition}[force Dimension]
  \label{def:physics:phyx-mini-0722:phyxminiproblems-problemphyxmini0722-forcedimension}
  \lean{PhyXMiniProblems.ProblemPhyXMini0722.forceDimension}
  The physical dimension of force, M L T⁻²
\end{definition}
\begin{proof}
  This is the declared model definition of force Dimension.
\end{proof}
\begin{definition}[Length Quantity]
  \label{def:physics:phyx-mini-0722:phyxminiproblems-problemphyxmini0722-lengthquantity}
  \lean{PhyXMiniProblems.ProblemPhyXMini0722.LengthQuantity}
  A nonnegative unit-independent physical length
\end{definition}
\begin{proof}
  This is the declared model definition of Length Quantity.
\end{proof}
\begin{definition}[Volume Quantity]
  \label{def:physics:phyx-mini-0722:phyxminiproblems-problemphyxmini0722-volumequantity}
  \lean{PhyXMiniProblems.ProblemPhyXMini0722.VolumeQuantity}
  \uses{def:physics:phyx-mini-0722:phyxminiproblems-problemphyxmini0722-volumedimension}
  A nonnegative unit-independent physical volume
\end{definition}
\begin{proof}
  This is the declared model definition of Volume Quantity.
\end{proof}
\begin{definition}[Mass Density Quantity]
  \label{def:physics:phyx-mini-0722:phyxminiproblems-problemphyxmini0722-massdensityquantity}
  \lean{PhyXMiniProblems.ProblemPhyXMini0722.MassDensityQuantity}
  \uses{def:physics:phyx-mini-0722:phyxminiproblems-problemphyxmini0722-massdensitydimension}
  A nonnegative unit-independent physical mass density
\end{definition}
\begin{proof}
  This is the declared model definition of Mass Density Quantity.
\end{proof}
\begin{definition}[Mass Quantity]
  \label{def:physics:phyx-mini-0722:phyxminiproblems-problemphyxmini0722-massquantity}
  \lean{PhyXMiniProblems.ProblemPhyXMini0722.MassQuantity}
  A nonnegative unit-independent physical mass
\end{definition}
\begin{proof}
  This is the declared model definition of Mass Quantity.
\end{proof}
\begin{definition}[Acceleration Quantity]
  \label{def:physics:phyx-mini-0722:phyxminiproblems-problemphyxmini0722-accelerationquantity}
  \lean{PhyXMiniProblems.ProblemPhyXMini0722.AccelerationQuantity}
  \uses{def:physics:phyx-mini-0722:phyxminiproblems-problemphyxmini0722-accelerationdimension}
  A nonnegative unit-independent physical acceleration magnitude
\end{definition}
\begin{proof}
  This is the declared model definition of Acceleration Quantity.
\end{proof}
\begin{definition}[Force Quantity]
  \label{def:physics:phyx-mini-0722:phyxminiproblems-problemphyxmini0722-forcequantity}
  \lean{PhyXMiniProblems.ProblemPhyXMini0722.ForceQuantity}
  \uses{def:physics:phyx-mini-0722:phyxminiproblems-problemphyxmini0722-forcedimension}
  A nonnegative unit-independent physical force magnitude
\end{definition}
\begin{proof}
  This is the declared model definition of Force Quantity.
\end{proof}
\begin{definition}[length Readout]
  \label{def:physics:phyx-mini-0722:phyxminiproblems-problemphyxmini0722-lengthreadout}
  \lean{PhyXMiniProblems.ProblemPhyXMini0722.lengthReadout}
  \uses{def:physics:phyx-mini-0722:phyxminiproblems-problemphyxmini0722-lengthquantity}
  Read a physical length in a selected length unit
\end{definition}
\begin{proof}
  This is the declared model definition of length Readout.
\end{proof}
\begin{definition}[volume Readout]
  \label{def:physics:phyx-mini-0722:phyxminiproblems-problemphyxmini0722-volumereadout}
  \lean{PhyXMiniProblems.ProblemPhyXMini0722.volumeReadout}
  \uses{def:physics:phyx-mini-0722:phyxminiproblems-problemphyxmini0722-volumequantity}
  Read a physical volume in the cube of a selected length unit
\end{definition}
\begin{proof}
  This is the declared model definition of volume Readout.
\end{proof}
\begin{definition}[density Readout]
  \label{def:physics:phyx-mini-0722:phyxminiproblems-problemphyxmini0722-densityreadout}
  \lean{PhyXMiniProblems.ProblemPhyXMini0722.densityReadout}
  \uses{def:physics:phyx-mini-0722:phyxminiproblems-problemphyxmini0722-massdensityquantity}
  Read density in a selected mass unit per selected length unit cubed
\end{definition}
\begin{proof}
  This is the declared model definition of density Readout.
\end{proof}
\begin{definition}[mass Readout]
  \label{def:physics:phyx-mini-0722:phyxminiproblems-problemphyxmini0722-massreadout}
  \lean{PhyXMiniProblems.ProblemPhyXMini0722.massReadout}
  \uses{def:physics:phyx-mini-0722:phyxminiproblems-problemphyxmini0722-massquantity}
  Read a physical mass in a selected mass unit
\end{definition}
\begin{proof}
  This is the declared model definition of mass Readout.
\end{proof}
\begin{definition}[acceleration Readout]
  \label{def:physics:phyx-mini-0722:phyxminiproblems-problemphyxmini0722-accelerationreadout}
  \lean{PhyXMiniProblems.ProblemPhyXMini0722.accelerationReadout}
  \uses{def:physics:phyx-mini-0722:phyxminiproblems-problemphyxmini0722-accelerationquantity}
  Read acceleration in coherent selected length and time units
\end{definition}
\begin{proof}
  This is the declared model definition of acceleration Readout.
\end{proof}
\begin{definition}[force Readout]
  \label{def:physics:phyx-mini-0722:phyxminiproblems-problemphyxmini0722-forcereadout}
  \lean{PhyXMiniProblems.ProblemPhyXMini0722.forceReadout}
  \uses{def:physics:phyx-mini-0722:phyxminiproblems-problemphyxmini0722-forcequantity}
  Read force in coherent selected mass, length, and time units
\end{definition}
\begin{proof}
  This is the declared model definition of force Readout.
\end{proof}
\begin{definition}[length In Centimeters]
  \label{def:physics:phyx-mini-0722:phyxminiproblems-problemphyxmini0722-lengthincentimeters}
  \lean{PhyXMiniProblems.ProblemPhyXMini0722.lengthInCentimeters}
  \uses{def:physics:phyx-mini-0722:phyxminiproblems-problemphyxmini0722-lengthquantity, def:physics:phyx-mini-0722:phyxminiproblems-problemphyxmini0722-lengthreadout}
  Centimeter readout of a physical length
\end{definition}
\begin{proof}
  This is the declared model definition of length In Centimeters.
\end{proof}
\begin{definition}[length In Meters]
  \label{def:physics:phyx-mini-0722:phyxminiproblems-problemphyxmini0722-lengthinmeters}
  \lean{PhyXMiniProblems.ProblemPhyXMini0722.lengthInMeters}
  \uses{def:physics:phyx-mini-0722:phyxminiproblems-problemphyxmini0722-lengthquantity, def:physics:phyx-mini-0722:phyxminiproblems-problemphyxmini0722-lengthreadout}
  Meter readout of a physical length
\end{definition}
\begin{proof}
  This is the declared model definition of length In Meters.
\end{proof}
\begin{definition}[volume In Cubic Meters]
  \label{def:physics:phyx-mini-0722:phyxminiproblems-problemphyxmini0722-volumeincubicmeters}
  \lean{PhyXMiniProblems.ProblemPhyXMini0722.volumeInCubicMeters}
  \uses{def:physics:phyx-mini-0722:phyxminiproblems-problemphyxmini0722-volumequantity, def:physics:phyx-mini-0722:phyxminiproblems-problemphyxmini0722-volumereadout}
  Cubic-meter readout of a physical volume
\end{definition}
\begin{proof}
  This is the declared model definition of volume In Cubic Meters.
\end{proof}
\begin{definition}[density In Kilograms Per Cubic Meter]
  \label{def:physics:phyx-mini-0722:phyxminiproblems-problemphyxmini0722-densityinkilogramspercubicmeter}
  \lean{PhyXMiniProblems.ProblemPhyXMini0722.densityInKilogramsPerCubicMeter}
  \uses{def:physics:phyx-mini-0722:phyxminiproblems-problemphyxmini0722-massdensityquantity, def:physics:phyx-mini-0722:phyxminiproblems-problemphyxmini0722-densityreadout}
  Kilogram-per-cubic-meter readout of a mass density
\end{definition}
\begin{proof}
  This is the declared model definition of density In Kilograms Per Cubic Meter.
\end{proof}
\begin{definition}[mass In Kilograms]
  \label{def:physics:phyx-mini-0722:phyxminiproblems-problemphyxmini0722-massinkilograms}
  \lean{PhyXMiniProblems.ProblemPhyXMini0722.massInKilograms}
  \uses{def:physics:phyx-mini-0722:phyxminiproblems-problemphyxmini0722-massquantity, def:physics:phyx-mini-0722:phyxminiproblems-problemphyxmini0722-massreadout}
  Kilogram readout of a physical mass
\end{definition}
\begin{proof}
  This is the declared model definition of mass In Kilograms.
\end{proof}
\begin{definition}[acceleration In Meters Per Second Squared]
  \label{def:physics:phyx-mini-0722:phyxminiproblems-problemphyxmini0722-accelerationinmeterspersecondsquared}
  \lean{PhyXMiniProblems.ProblemPhyXMini0722.accelerationInMetersPerSecondSquared}
  \uses{def:physics:phyx-mini-0722:phyxminiproblems-problemphyxmini0722-accelerationquantity, def:physics:phyx-mini-0722:phyxminiproblems-problemphyxmini0722-accelerationreadout}
  Meter-per-second-squared readout of an acceleration magnitude
\end{definition}
\begin{proof}
  This is the declared model definition of acceleration In Meters Per Second Squared.
\end{proof}
\begin{definition}[force In Newtons]
  \label{def:physics:phyx-mini-0722:phyxminiproblems-problemphyxmini0722-forceinnewtons}
  \lean{PhyXMiniProblems.ProblemPhyXMini0722.forceInNewtons}
  \uses{def:physics:phyx-mini-0722:phyxminiproblems-problemphyxmini0722-forcequantity, def:physics:phyx-mini-0722:phyxminiproblems-problemphyxmini0722-forcereadout}
  Newton readout of a physical force magnitude
\end{definition}
\begin{proof}
  This is the declared model definition of force In Newtons.
\end{proof}
\begin{definition}[Block Material]
  \label{def:physics:phyx-mini-0722:phyxminiproblems-problemphyxmini0722-blockmaterial}
  \lean{PhyXMiniProblems.ProblemPhyXMini0722.BlockMaterial}
  The material of the submerged block
\end{definition}
\begin{proof}
  This is the declared model definition of Block Material.
\end{proof}
\begin{definition}[Fluid Kind]
  \label{def:physics:phyx-mini-0722:phyxminiproblems-problemphyxmini0722-fluidkind}
  \lean{PhyXMiniProblems.ProblemPhyXMini0722.FluidKind}
  The liquid surrounding the block
\end{definition}
\begin{proof}
  This is the declared model definition of Fluid Kind.
\end{proof}
\begin{definition}[Immersion State]
  \label{def:physics:phyx-mini-0722:phyxminiproblems-problemphyxmini0722-immersionstate}
  \lean{PhyXMiniProblems.ProblemPhyXMini0722.ImmersionState}
  Whether the block displaces its full geometric volume
\end{definition}
\begin{proof}
  This is the declared model definition of Immersion State.
\end{proof}
\begin{definition}[String Attachment]
  \label{def:physics:phyx-mini-0722:phyxminiproblems-problemphyxmini0722-stringattachment}
  \lean{PhyXMiniProblems.ProblemPhyXMini0722.StringAttachment}
  The anchoring geometry of the string
\end{definition}
\begin{proof}
  This is the declared model definition of String Attachment.
\end{proof}
\begin{definition}[Motion State]
  \label{def:physics:phyx-mini-0722:phyxminiproblems-problemphyxmini0722-motionstate}
  \lean{PhyXMiniProblems.ProblemPhyXMini0722.MotionState}
  The translational state relevant to the force balance
\end{definition}
\begin{proof}
  This is the declared model definition of Motion State.
\end{proof}
\begin{definition}[Vertical Direction]
  \label{def:physics:phyx-mini-0722:phyxminiproblems-problemphyxmini0722-verticaldirection}
  \lean{PhyXMiniProblems.ProblemPhyXMini0722.VerticalDirection}
  Vertical directions used by the force arrows in image 722
\end{definition}
\begin{proof}
  This is the declared model definition of Vertical Direction.
\end{proof}
\begin{definition}[Figure Object]
  \label{def:physics:phyx-mini-0722:phyxminiproblems-problemphyxmini0722-figureobject}
  \lean{PhyXMiniProblems.ProblemPhyXMini0722.FigureObject}
  Physical objects and graphical elements visible in image 722
\end{definition}
\begin{proof}
  This is the declared model definition of Figure Object.
\end{proof}
\begin{definition}[Figure Label]
  \label{def:physics:phyx-mini-0722:phyxminiproblems-problemphyxmini0722-figurelabel}
  \lean{PhyXMiniProblems.ProblemPhyXMini0722.FigureLabel}
  Text and force-symbol labels visible in image 722
\end{definition}
\begin{proof}
  This is the declared model definition of Figure Label.
\end{proof}
\begin{definition}[Supplied Submerged Block Figure]
  \label{def:physics:phyx-mini-0722:phyxminiproblems-problemphyxmini0722-suppliedsubmergedblockfigure}
  \lean{PhyXMiniProblems.ProblemPhyXMini0722.SuppliedSubmergedBlockFigure}
  \uses{def:physics:phyx-mini-0722:phyxminiproblems-problemphyxmini0722-verticaldirection, def:physics:phyx-mini-0722:phyxminiproblems-problemphyxmini0722-figureobject, def:physics:phyx-mini-0722:phyxminiproblems-problemphyxmini0722-figurelabel}
  Qualitative evidence transcribed from the primary raster. The image contains force labels and directions but no numerical value of the requested tension
\end{definition}
\begin{proof}
  This is the declared model definition of Supplied Submerged Block Figure.
\end{proof}
\begin{definition}[Submerged Wood Block Setup]
  \label{def:physics:phyx-mini-0722:phyxminiproblems-problemphyxmini0722-submergedwoodblocksetup}
  \lean{PhyXMiniProblems.ProblemPhyXMini0722.SubmergedWoodBlockSetup}
  \uses{def:physics:phyx-mini-0722:phyxminiproblems-problemphyxmini0722-lengthquantity, def:physics:phyx-mini-0722:phyxminiproblems-problemphyxmini0722-volumequantity, def:physics:phyx-mini-0722:phyxminiproblems-problemphyxmini0722-massdensityquantity, def:physics:phyx-mini-0722:phyxminiproblems-problemphyxmini0722-massquantity, def:physics:phyx-mini-0722:phyxminiproblems-problemphyxmini0722-accelerationquantity, def:physics:phyx-mini-0722:phyxminiproblems-problemphyxmini0722-forcequantity, def:physics:phyx-mini-0722:phyxminiproblems-problemphyxmini0722-blockmaterial, def:physics:phyx-mini-0722:phyxminiproblems-problemphyxmini0722-fluidkind, def:physics:phyx-mini-0722:phyxminiproblems-problemphyxmini0722-immersionstate, def:physics:phyx-mini-0722:phyxminiproblems-problemphyxmini0722-stringattachment, def:physics:phyx-mini-0722:phyxminiproblems-problemphyxmini0722-motionstate, def:physics:phyx-mini-0722:phyxminiproblems-problemphyxmini0722-suppliedsubmergedblockfigure}
  Independent quantities and observables of the submerged-block experiment. None of the force fields is defined from an answer choice
\end{definition}
\begin{proof}
  This is the declared model definition of Submerged Wood Block Setup.
\end{proof}
\begin{definition}[Matches Submerged Wood Block Scenario]
  \label{def:physics:phyx-mini-0722:phyxminiproblems-problemphyxmini0722-matchessubmergedwoodblockscenario}
  \lean{PhyXMiniProblems.ProblemPhyXMini0722.MatchesSubmergedWoodBlockScenario}
  \uses{def:physics:phyx-mini-0722:phyxminiproblems-problemphyxmini0722-submergedwoodblocksetup}
  Categorical facts stated by the prose and shown in the figure
\end{definition}
\begin{proof}
  This is the declared model definition of Matches Submerged Wood Block Scenario.
\end{proof}
\begin{definition}[Matches Supplied Submerged Block Figure]
  \label{def:physics:phyx-mini-0722:phyxminiproblems-problemphyxmini0722-matchessuppliedsubmergedblockfigure}
  \lean{PhyXMiniProblems.ProblemPhyXMini0722.MatchesSuppliedSubmergedBlockFigure}
  \uses{def:physics:phyx-mini-0722:phyxminiproblems-problemphyxmini0722-submergedwoodblocksetup}
  Objects, labels, and force directions read from the primary image
\end{definition}
\begin{proof}
  This is the declared model definition of Matches Supplied Submerged Block Figure.
\end{proof}
\begin{definition}[Matches Submerged Block Problem Readouts]
  \label{def:physics:phyx-mini-0722:phyxminiproblems-problemphyxmini0722-matchessubmergedblockproblemreadouts}
  \lean{PhyXMiniProblems.ProblemPhyXMini0722.MatchesSubmergedBlockProblemReadouts}
  \uses{def:physics:phyx-mini-0722:phyxminiproblems-problemphyxmini0722-lengthincentimeters, def:physics:phyx-mini-0722:phyxminiproblems-problemphyxmini0722-densityinkilogramspercubicmeter, def:physics:phyx-mini-0722:phyxminiproblems-problemphyxmini0722-submergedwoodblocksetup}
  Numerical readouts explicitly given in the problem. Keeping all three side lengths preserves the 10 cm × 10 cm × 10 cm geometry
\end{definition}
\begin{proof}
  This is the declared model definition of Matches Submerged Block Problem Readouts.
\end{proof}
\begin{definition}[Uses Standard Water And Gravity]
  \label{def:physics:phyx-mini-0722:phyxminiproblems-problemphyxmini0722-usesstandardwaterandgravity}
  \lean{PhyXMiniProblems.ProblemPhyXMini0722.UsesStandardWaterAndGravity}
  \uses{def:physics:phyx-mini-0722:phyxminiproblems-problemphyxmini0722-densityinkilogramspercubicmeter, def:physics:phyx-mini-0722:phyxminiproblems-problemphyxmini0722-accelerationinmeterspersecondsquared, def:physics:phyx-mini-0722:phyxminiproblems-problemphyxmini0722-submergedwoodblocksetup}
  Textbook calibrations implicit in the recorded answer: freshwater density is 1000 kg/m³ and near-Earth gravitational acceleration is 9.8 m/s². They are kept separate from the values printed in the question
\end{definition}
\begin{proof}
  This is the declared model definition of Uses Standard Water And Gravity.
\end{proof}
\begin{definition}[Has Physical Submerged Block Parameters]
  \label{def:physics:phyx-mini-0722:phyxminiproblems-problemphyxmini0722-hasphysicalsubmergedblockparameters}
  \lean{PhyXMiniProblems.ProblemPhyXMini0722.HasPhysicalSubmergedBlockParameters}
  \uses{def:physics:phyx-mini-0722:phyxminiproblems-problemphyxmini0722-lengthinmeters, def:physics:phyx-mini-0722:phyxminiproblems-problemphyxmini0722-volumeincubicmeters, def:physics:phyx-mini-0722:phyxminiproblems-problemphyxmini0722-densityinkilogramspercubicmeter, def:physics:phyx-mini-0722:phyxminiproblems-problemphyxmini0722-massinkilograms, def:physics:phyx-mini-0722:phyxminiproblems-problemphyxmini0722-accelerationinmeterspersecondsquared, def:physics:phyx-mini-0722:phyxminiproblems-problemphyxmini0722-submergedwoodblocksetup}
  Positivity and nondegeneracy conditions for the physical setup
\end{definition}
\begin{proof}
  This is the declared model definition of Has Physical Submerged Block Parameters.
\end{proof}
\begin{definition}[Satisfies Cube And Displacement Laws]
  \label{def:physics:phyx-mini-0722:phyxminiproblems-problemphyxmini0722-satisfiescubeanddisplacementlaws}
  \lean{PhyXMiniProblems.ProblemPhyXMini0722.SatisfiesCubeAndDisplacementLaws}
  \uses{def:physics:phyx-mini-0722:phyxminiproblems-problemphyxmini0722-lengthreadout, def:physics:phyx-mini-0722:phyxminiproblems-problemphyxmini0722-volumereadout, def:physics:phyx-mini-0722:phyxminiproblems-problemphyxmini0722-submergedwoodblocksetup}
  Cube-volume and full-submersion displacement laws, stated in every coherent length unit. These laws do not mention any force result
\end{definition}
\begin{proof}
  This is the declared model definition of Satisfies Cube And Displacement Laws.
\end{proof}
\begin{definition}[Satisfies Submerged Block Statics Laws]
  \label{def:physics:phyx-mini-0722:phyxminiproblems-problemphyxmini0722-satisfiessubmergedblockstaticslaws}
  \lean{PhyXMiniProblems.ProblemPhyXMini0722.SatisfiesSubmergedBlockStaticsLaws}
  \uses{def:physics:phyx-mini-0722:phyxminiproblems-problemphyxmini0722-volumereadout, def:physics:phyx-mini-0722:phyxminiproblems-problemphyxmini0722-densityreadout, def:physics:phyx-mini-0722:phyxminiproblems-problemphyxmini0722-massreadout, def:physics:phyx-mini-0722:phyxminiproblems-problemphyxmini0722-accelerationreadout, def:physics:phyx-mini-0722:phyxminiproblems-problemphyxmini0722-forcereadout, def:physics:phyx-mini-0722:phyxminiproblems-problemphyxmini0722-submergedwoodblocksetup}
  The governing hydrostatic and statics laws: mass from density and volume, weight from mass and gravity, Archimedes' buoyant-force law, and the vertical force balance for a block held at rest. None contains the requested tension or an answer-choice number
\end{definition}
\begin{proof}
  This is the declared model definition of Satisfies Submerged Block Statics Laws.
\end{proof}
\begin{definition}[Answer Choice]
  \label{def:physics:phyx-mini-0722:phyxminiproblems-problemphyxmini0722-answerchoice}
  \lean{PhyXMiniProblems.ProblemPhyXMini0722.AnswerChoice}
  Labels of the four force alternatives printed with the problem
\end{definition}
\begin{proof}
  This is the declared model definition of Answer Choice.
\end{proof}
\begin{definition}[displayed Tension In Newtons]
  \label{def:physics:phyx-mini-0722:phyxminiproblems-problemphyxmini0722-displayedtensioninnewtons}
  \lean{PhyXMiniProblems.ProblemPhyXMini0722.displayedTensionInNewtons}
  \uses{def:physics:phyx-mini-0722:phyxminiproblems-problemphyxmini0722-answerchoice}
  Tension in newtons printed beside each answer label
\end{definition}
\begin{proof}
  This is the declared model definition of displayed Tension In Newtons.
\end{proof}
\begin{definition}[recorded Dataset Answer]
  \label{def:physics:phyx-mini-0722:phyxminiproblems-problemphyxmini0722-recordeddatasetanswer}
  \lean{PhyXMiniProblems.ProblemPhyXMini0722.recordedDatasetAnswer}
  \uses{def:physics:phyx-mini-0722:phyxminiproblems-problemphyxmini0722-answerchoice}
  Dataset answer metadata, deliberately not used as a theorem premise
\end{definition}
\begin{proof}
  This is the declared model definition of recorded Dataset Answer.
\end{proof}
\begin{definition}[Rounds To Nearest Tenth Newton]
  \label{def:physics:phyx-mini-0722:phyxminiproblems-problemphyxmini0722-roundstonearesttenthnewton}
  \lean{PhyXMiniProblems.ProblemPhyXMini0722.RoundsToNearestTenthNewton}
  \uses{def:physics:phyx-mini-0722:phyxminiproblems-problemphyxmini0722-forcequantity, def:physics:phyx-mini-0722:phyxminiproblems-problemphyxmini0722-forceinnewtons}
  Agreement with a force displayed to the nearest tenth of a newton
\end{definition}
\begin{proof}
  This is the declared model definition of Rounds To Nearest Tenth Newton.
\end{proof}
\begin{definition}[Is Nearest Displayed Tension]
  \label{def:physics:phyx-mini-0722:phyxminiproblems-problemphyxmini0722-isnearestdisplayedtension}
  \lean{PhyXMiniProblems.ProblemPhyXMini0722.IsNearestDisplayedTension}
  \uses{def:physics:phyx-mini-0722:phyxminiproblems-problemphyxmini0722-forceinnewtons, def:physics:phyx-mini-0722:phyxminiproblems-problemphyxmini0722-submergedwoodblocksetup, def:physics:phyx-mini-0722:phyxminiproblems-problemphyxmini0722-answerchoice, def:physics:phyx-mini-0722:phyxminiproblems-problemphyxmini0722-displayedtensioninnewtons}
  A displayed choice is at least as close as every alternative
\end{definition}
\begin{proof}
  This is the declared model definition of Is Nearest Displayed Tension.
\end{proof}
% --- Archon declaration topology end ---
% --- Archon physics formalization source end ---
